\scanpage{67}
% archive sheet number: 69

et considérons ``l'espace de Teichmüller à bord'' correspondant \deleted{$\unsure{T_{g,\nu}|T_{g,r}}$} $TB_{g,\nu}$, qui est \deleted{isom.} \inserted{homéo.} à $\mathbb R^{6g-6+\unsure{3|2}\nu}$ ou mieux, à $\mathbb C^{3g-3+\nu}\times\mathbb R^\nu$. On a une application
\diagram{\begin{tikzcd}[ampersand replacement=\&,column sep=large,row sep=0pt]
TB_{g,\nu}\arrow[r] \& T_{g\unsure{\,|,}\nu}\\
\simeq\mathbb C^{3g-3+\nu}\times\mathbb R^\nu \& \mathbb C^{3g-3+\nu}\deleted{\times\mathbb R^\nu}
\end{tikzcd}}
de bouchage des trous, je présume qu'elle est fibrante \deleted{\unsure{de}} façon, plus précise, que \deleted{deux pts de} \inserted{c'est la projection} canonique. \deleted{\unsure{de|et}} En effet, cela doit provenir du fait que pour une sphère \inserted{\unsure{triforée|tri-trouée|holomorphe}} à bord de type $(0,3)$, \deleted{dont} quand le modèle combinatoire (\unsure{triangle|triangles} combinatoire) est fixé, et qu'on varie seulement \deleted{le} un, deux ou trois des paramètres ``longueurs'' correspondants aux trois trous, la sphère ``à bord et ponctuée'' (mono, bi ou triponctuée suivant les cas) \deleted{est la} \inserted{\unsure{associée|congruente}}
